\item \textbf{Datación por Carbono-14 de celulosa antigua.}
Un espécimen vivo en equilibrio con la atmósfera posee una razón de un átomo de $^{14}\text{C}$ ($T_{1/2} = 5730\text{ años}$) por cada $7.7 \times 10^{11}$ átomos de carbono estable $^{12}\text{C}$. Una muestra arqueológica de madera antigua (celulosa, $\text{C}_{12}\text{H}_{22}\text{O}_{11}$) contiene exactamente $21.0\text{ mg}$ de carbono. Al introducirla en un contador beta con una eficiencia de detección del $88.0\%$, se registran $837$ desintegraciones en el transcurso de una semana.
\begin{enumerate}
    \item Encuentre el número total de átomos de carbono en la muestra.
    \item Encuentre la cantidad de átomos de $^{14}\text{C}$ presentes originalmente en la muestra.
    \item Determine la constante de desintegración $\lambda$ del $^{14}\text{C}$ en $\text{s}^{-1}$.
    \item Halle el número de desintegraciones semanales iniciales al morir el espécimen.
    \item Determine la tasa corregida de desintegraciones por semana actual medida en el detector.
    \item A partir de lo anterior, determine con precisión la antigüedad en años de la muestra arqueológica.
\end{enumerate}